\chapter{Conditional Expectation}

Conditional expectation is arguably the most important concept in mathematical finance, as prices of fair assets are modeled as conditional expectations of their future payoffs under a specific measure (martingales).

\section{Conditional Expectation in \texorpdfstring{$L^1$}{L1}}

The elementary definition $\mathbb{E}[X|Y=y] = \sum x \mathbb{P}(X=x|Y=y)$ works for discrete variables but fails for continuous ones because $\mathbb{P}(Y=y) = 0$. The rigorous definition relies on $\sigma$-algebras and measurability.

\begin{theorem}[Existence and Uniqueness]
Let $(\Omega, \mathcal{F}, \mathbb{P})$ be a probability space and let $X \in L^1(\Omega, \mathcal{F}, \mathbb{P})$. Let $\mathcal{G} \subseteq \mathcal{F}$ be a sub-$\sigma$-algebra (representing ``partial information'').

There exists a \textbf{unique} random variable $Z$, denoted by $\mathbb{E}[X|\mathcal{G}]$, such that:
\begin{enumerate}
    \item \textbf{Measurability}: $Z$ is $\mathcal{G}$-measurable.
    \item \textbf{Partial Averaging Property}: For every bounded $\mathcal{G}$-measurable random variable $H$ (or equivalently for all $A \in \mathcal{G}$):
    \[
    \mathbb{E}[Z H] = \mathbb{E}[X H] \quad \left( \text{i.e., } \int_A Z \, d\mathbb{P} = \int_A X \, d\mathbb{P} \right).
    \]
\end{enumerate}
This $Z$ is called the \textbf{conditional expectation} of $X$ given $\mathcal{G}$.
\end{theorem}
\textit{Uniqueness holds $\mathbb{P}$-almost surely.}

\paragraph{Intuition: What does \texorpdfstring{$\mathcal{G}$}{G}-measurable mean?}
It means that the value of $Z$ is fully determined by the information in $\mathcal{G}$.
\begin{itemize}
    \item If $\mathcal{G}$ represents our knowledge at time $t$, then $Z$ is a value we can \textbf{observe right now}.
    \item If $\mathcal{G} = \sigma(Y)$, then $Z$ can be written as a deterministic function of $Y$: $Z = f(Y)$.
    \item In simple terms: "Given $\mathcal{G}$, there is no more randomness in $Z$".
\end{itemize}

\paragraph{Intuition: The "Partial Averaging" Concept}
Imagine $\mathcal{G}$ is generated by a partition of $\Omega$ into disjoint sets $A_1, A_2, \dots$.
\begin{itemize}
    \item On each set $A_k$, we cannot distinguish the specific outcomes $\omega$. The available information is just "we are in $A_k$".
    \item Therefore, the best guess $Z = \mathbb{E}[X|\mathcal{G}]$ must be \textbf{constant} on $A_k$.
    \item To satisfy the property $\int_{A_k} Z d\mathbb{P} = \int_{A_k} X d\mathbb{P}$, this constant must be the \textbf{average value} of $X$ on $A_k$:
    \[
    Z(\omega) = \frac{1}{\mathbb{P}(A_k)} \int_{A_k} X(\omega') d\mathbb{P}(\omega') \quad \text{for } \omega \in A_k.
    \]
    \item In general $\sigma$-algebras (not just partitions), $Z$ is a "smoothed" version of $X$ that preserves the local averages over all $\mathcal{G}$-measurable sets.
\end{itemize}

\paragraph{Example: Die Roll Partition}
Consider rolling a fair die $\Omega = \{1, 2, 3, 4, 5, 6\}$.
\begin{itemize}
    \item Let the information be "Is the number Even or Odd?". This generates a partition of $\Omega$:
    \[
    A_1 = \{1, 3, 5\} \text{ (Odd)}, \quad A_2 = \{2, 4, 6\} \text{ (Even)}.
    \]
    \item These sets are \textbf{disjoint} ($A_1 \cap A_2 = \emptyset$) and their union is the whole space $\Omega$.
    \item If we know we are in $A_2$ (Even), our best guess for the outcome is the average of $\{2, 4, 6\}$, which is $4$.
    \item So $\mathbb{E}[X | \text{"Even"}] = 4$.
\end{itemize}

\begin{remark}[Generated $\sigma$-algebra]
If $\mathcal{G} = \sigma(Y)$, this is the \textbf{$\sigma$-algebra generated by $Y$}.
\begin{itemize}
    \item It represents the information provided by observing the random variable $Y$.
    \item Formally, $\sigma(Y) = \{ Y^{-1}(B) : B \in \mathcal{B}(\mathbb{R}) \}$.
    \item Conditioning on $Y$ means conditioning on the information "we know the value of $Y$".
\end{itemize}
We write $\mathbb{E}[X|Y]$ instead of $\mathbb{E}[X|\sigma(Y)]$ for brevity.
\end{remark}

\section{Key Properties}
Let $X, Y \in L^1$ and $\mathcal{G}, \mathcal{H}$ be sub-$\sigma$-algebras.

\begin{proposition}
The conditional expectation satisfies:
\begin{enumerate}
    \item \textbf{Linearity}: $\mathbb{E}[aX + bY | \mathcal{G}] = a\mathbb{E}[X|\mathcal{G}] + b\mathbb{E}[Y|\mathcal{G}]$.
    \item \textbf{Positivity}: If $X \geq 0$ a.s., then $\mathbb{E}[X|\mathcal{G}] \geq 0$ a.s.
    \item \textbf{Taking out what is known}: If $Y$ is $\mathcal{G}$-measurable and $XY \in L^1$, then:
    \[
    \mathbb{E}[XY|\mathcal{G}] = Y \mathbb{E}[X|\mathcal{G}].
    \]
    \item \textbf{Tower Property} (Iterated Conditioning): If $\mathcal{H} \subset \mathcal{G}$ (less information in $\mathcal{H}$), then:
    \[
    \mathbb{E}[ \mathbb{E}[X|\mathcal{G}] | \mathcal{H} ] = \mathbb{E}[X|\mathcal{H}].
    \]
    Evaluating the best guess of a best guess yields the coarser guess.
    \item \textbf{Independence}: If $X$ is independent of $\mathcal{G}$, then $\mathbb{E}[X|\mathcal{G}] = \mathbb{E}[X]$.
    \item \textbf{Full Information}: If $X$ is $\mathcal{G}$-measurable, then $\mathbb{E}[X|\mathcal{G}] = X$.
    \item \textbf{Jensen's Inequality}: If $\phi: \mathbb{R} \to \mathbb{R}$ is convex, then $\phi(\mathbb{E}[X|\mathcal{G}]) \leq \mathbb{E}[\phi(X)|\mathcal{G}]$.
\end{enumerate}
\end{proposition}

\section{Geometric Interpretation in \texorpdfstring{$L^2$}{L2}}

If we restrict ourselves to $L^2(\Omega, \mathcal{F}, \mathbb{P})$, the space is a Hilbert space. The set of $\mathcal{G}$-measurable random variables in $L^2$, denoted $L^2(\mathcal{G})$, is a closed linear subspace of $L^2(\mathcal{F})$.

\begin{theorem}[Orthogonal Projection]
For $X \in L^2(\mathcal{F})$, the conditional expectation $\mathbb{E}[X|\mathcal{G}]$ is the \textbf{orthogonal projection} of $X$ onto the subspace $L^2(\mathcal{G})$.
\[
\mathbb{E}[X|\mathcal{G}] = \text{argmin}_{Z \in L^2(\mathcal{G})} \| X - Z \|_2 = \text{argmin}_{Z \in L^2(\mathcal{G})} \mathbb{E}[(X-Z)^2].
\]
Using Hilbert space geometry, this means the error $X - \mathbb{E}[X|\mathcal{G}]$ is orthogonal to every vector in the subspace $L^2(\mathcal{G})$:
\[
\langle X - \mathbb{E}[X|\mathcal{G}], Z \rangle = 0 \quad \forall Z \in L^2(\mathcal{G}).
\]
\end{theorem}
This interprets the conditional expectation as the ``best Mean-Square estimator'' of $X$ given the information $\mathcal{G}$.

\subsection{Example: Best Constant Approximation}
Consider the trivial $\sigma$-algebra $\mathcal{G} = \{ \emptyset, \Omega \}$ (no information).
\begin{itemize}
    \item The space $L^2(\mathcal{G})$ consists only of \textbf{constant random variables} $Z(\omega) = c$.
    \item The orthogonal projection problem becomes finding the best constant $c$ to approximate $X$:
    \[
    \min_{c \in \mathbb{R}} \mathbb{E}[(X - c)^2].
    \]
    \item By simple calculus (differentiating with respect to $c$), the minimum is achieved at $c = \mathbb{E}[X]$.
    \item \textbf{Geometric Interpretation}: The expectation $\mathbb{E}[X]$ is the orthogonal projection of the vector $X$ onto the line of constants.
    \item The Pythagorean theorem holds: $\|X\|^2 = \|\mathbb{E}[X]\|^2 + \|X - \mathbb{E}[X]\|^2$, which corresponds to $\mathbb{E}[X^2] = (\mathbb{E}[X])^2 + \text{Var}(X)$.
\end{itemize}

\section{Gaussian Regression}
A powerful application of the $L^2$ projection view is in Gaussian vectors.

\begin{theorem}
Let $(Y_1, \dots, Y_n, X)$ be a Gaussian vector in $L^2$. We want to estimate $X$ given observations $Y = (Y_1, \dots, Y_n)^T$.
The conditional expectation is linear in $Y$:
\[
\mathbb{E}[X | Y_1, \dots, Y_n] = \mathbb{E}[X] + \lambda^T (Y - \mathbb{E}[Y])
\]
where $\lambda = \Sigma_Y^{-1} \text{Cov}(X, Y)$. Here, $\Sigma_Y$ is the covariance matrix of $Y$, and $\text{Cov}(X, Y)$ is the vector of covariances between $X$ and each $Y_i$.
\end{theorem}

\subsection{Example: Bivariate Normal Case}
Let $(X, Y)$ be a bivariate Gaussian vector with means $\mu_X, \mu_Y$, variances $\sigma_X^2, \sigma_Y^2$, and correlation $\rho$.
\begin{itemize}
    \item We want to estimate $X$ given $Y$.
    \item The covariance is $\text{Cov}(X, Y) = \rho \sigma_X \sigma_Y$.
    \item The variance of $Y$ is $\Sigma_Y = \sigma_Y^2$.
    \item Following the theorem, $\lambda = \frac{\rho \sigma_X \sigma_Y}{\sigma_Y^2} = \rho \frac{\sigma_X}{\sigma_Y}$.
\end{itemize}
Thus, the best estimator is linear:
\[
\mathbb{E}[X|Y] = \mu_X + \rho \frac{\sigma_X}{\sigma_Y} (Y - \mu_Y).
\]
\begin{itemize}
    \item If $\rho=0$ (independent), $\mathbb{E}[X|Y] = \mu_X$ (observation $Y$ usesless).
    \item If $\rho=1$ (perfect correlation), $\mathbb{E}[X|Y]$ perfectly reconstructs $X$ (up to scaling).
\end{itemize}
Moreover, the residual error $X - \mathbb{E}[X|Y]$ is independent of $\sigma(Y)$ (a specific property of Gaussianity; essentially orthogonality implies independence for Gaussians).
